\section{Verification, editorial reconciliation, and integration}
\label{sec:audit}

\subsection{What the computations establish}

The analytic theorems are proved in the text. Exact computation enters
at three clearly identified places: enclosure of the Gaussian maximum,
endpoint signs for the two Lerch branch-shape decisions, and separation
of a finite formal relation span. The distribution examples and inverse
coefficient checks are independent consistency tests of general proofs.

\begin{longtable}{@{}L{0.24\textwidth}L{0.37\textwidth}L{0.32\textwidth}@{}}
\toprule
Check & Construction & Logical role\\
\midrule
\endhead
Gaussian maximum & Outward rational intervals, finite Euler transforms,
analytic derivative tails & Certifies the printed root and maximum
intervals once strict log-concavity is proved.\\[4pt]
Inverse expansion & Exact symbolic generalized series, with exponents
encoded by rational bases & Substitution cancels every coefficient below
exponent two; the all-order theorem has an analytic proof.\\[4pt]
Cyclic quotients & Exact Smith forms of the original orbit relation
matrices, without a normal-form implementation & 52 scalar and finite-jet
cases check the theorem; three ramified examples test its boundary.\\[4pt]
Lerch branch shapes & Exact Euler--Maclaurin intervals on entire rational
root brackets & Certifies the signs deciding the order-two and
order-three global branch shapes.\\[4pt]
$S_6$ separation & Integer functional, independent row reconstruction
and complete-family enumeration & Proves nonmembership in the stated
5,131-row rational span.\\[4pt]
Plots and quadrature & Floating-point evaluation of independent
integrals and asymptotic formulas & Diagnostic evidence only; no claimed
certificate follows from a plot or an unenclosed decimal.\\
\bottomrule
\end{longtable}

The nine Gaussian comparison points evaluate the original Euler
transform and an independent beta-mixture quadrature. Their largest
observed discrepancy is below $6\cdot10^{-12}$ with the supplied
quadrature settings. This is a diagnostic discrepancy, not a proved
bound on the quadrature error. Four inverse examples compare the exact
coefficient formula with numerical solutions. The Lerch plot retains
its sampled values and analytic omitted-function estimates; floating
rounding is explicitly unenclosed. In particular, the diagnostic
coordinates of the order-two minimum are not a certified isolating box.
Its existence, uniqueness, and nondegeneracy are theorems independently
of those coordinates.

The programs and receipts are included in the archive. Running
\path{code/replay_all.py} from the package root regenerates all exact
receipts. Optional plotting dependencies are separated from the exact
arithmetic dependencies. No network access or numerical
integer-relation search is needed for the replay. The distribution
checks use SymPy's exact integer Smith algorithm; the Gaussian, Lerch,
and $S_6$ certificate arithmetic uses the Python standard library.

This is stronger than reporting successful high-precision fitting,
but it remains ordinary mathematical proof with auditable programs.
Formal verification in a proof assistant is a separate project. For
that project, the integer separator is an especially small logical
target: the checker need only reconstruct finite vectors and evaluate
integer dot products.

\subsection{Concrete proposed status corrections}

The reviewed material did not reveal a new counterexample to an
already established upstream theorem in the portions used here.
The useful corrections concern cross-report status and the scope of
claims made while integrating results.

\paragraph{Promote the critical constant conjecture.}
In \cite{realorder}, replace the status of Conjecture
\texttt{conj:constant} by a theorem, with the proof through
Theorems~\ref{thm:allocation} and \ref{thm:criticalconstant}.
Retain the analytic/Abel convention on $a+b=1$ and the endpoint atom.
The constant is a supremum over the open parameter interval, approached
as $a\downarrow0$; there is no admissible maximizing pair on the line.

\paragraph{Reconcile the first-order Lerch minimum.}
The formal-reductions report \cite{formalreductions} still lists
uniqueness of the smaller $n=2,k=1$ branch minimum as Conjecture
\texttt{research:conj:minimum}. The incoming boundary report
\cite{lershape:Boundary} proves that uniqueness. A combined manuscript
should cite the latter result and remove the obsolete open status.
The present contribution is the all-$k$ dichotomy, the endpoint
criterion, the order-two minimum, and the order-three decreasing
case. These should not be described as the first proof of the
already settled $k=1$ case.

\paragraph{Keep the $S_4$ and $S_6$ statuses distinct.}
The current manuscript has a proof chapter for $S_4$ and an explicitly
conjectural $S_6$ reduction in
\path{chapters/04-cyclotomic-quotients.tex}. The exact separator in
Section~\ref{s6bounded:sec} is a theorem about one finite presentation.
It supplies neither a proof nor a disproof of the $S_6$ evaluation.
Its enlarged row family should replace a repeated search in the same
smaller ansatz, rather than change the conjecture's numerical status.

\paragraph{Separate two different Euler constants.}
The global constant $C_*$ controls $-2g_{a,b}$ over all positive orders.
The critical constant $\pi/4+(\log2)/2$ controls every normalized Euler
tail on $a+b=1$. The exact $(a,b)=(1,1)$ calculation in
Section~\ref{sec:secant} proves that normalized errors need not decrease
off the critical line. This is a warning against an additional claim,
not a counterexample to either theorem stated here or to an identified
upstream theorem.

\paragraph{State the cyclic hypotheses where the formula is used.}
The restrictions $(m,d)=1$, $(u-1,d)=1$, and
$\ell\nmid\varphi(m)$ are part of the theorem. They should accompany
any extracted formula or summary. The level-nine computation shows
why a partially fixed primary factor is a distinct local problem.
Likewise, the term ``rank'' should specify whether it is a free
abelian rank, a formal rational dimension, or a rank after reduction
modulo a prime. No one of these determines the others without a proof.

\subsection{Suggested repository placement}

The complete report can first be placed under
\path{Analysis/Polylogarithms/docs/reports/extremal-bounds-cyclic-descent/}.
The integration guide supplies a compact addendum and a dependency
map. The natural manuscript neighborhoods are:
\begin{itemize}
\item Gaussian allocation, the critical Euler theorem, and the secant
identity next to the signed-kernel material and the real-order report;
\item the cyclic theorem and finite jets next to
\path{08-distribution.tex} and \path{08-conductor-jets.tex};
\item the Lerch branch theorems next to
\path{09-zero-geometry.tex} and the incoming boundary analysis;
\item the finite $S_6$ obstruction next to
\path{04-cyclotomic-quotients.tex} and the established $S_4$ proof.
\end{itemize}
The manuscript need not absorb the entire report at once. Its theorem
statements and their proof dependencies can be integrated by topic,
while the exact checkers and complete provenance remain with the report.
No repository mutation is part of this deliverable.

\section{Further research questions and proposed directions}
\label{sec:research}

The questions below are intentionally differentiated. Some preserve
existing conjectures, some ask for stronger theorems suggested by the
present proofs, and some concern a concrete missing computation.
No conjecture in this section is used earlier.

\subsection{Sharp bounds for the entire Euler scheme}

\paragraph{1. Optimize every normalized remainder.}
For $N\ge1$, determine
\[
 M_N=\sup_{a,b>0,\ a+b\ge1}2^N(E_N-g_{a,b}),
 \qquad M=\sup_{N\ge1}M_N.
\]
We have determined $M_1=C_*$. On the critical line the corresponding
all-$N$ supremum is $\pi/4+(\log2)/2$. The positive secant identity
reduces the general problem to a beta average of divided differences
of $H_N$. The main obstacle is concrete: $H_1$ is concave on
$[1,\infty)$, whereas $H_2$ is convex near one. A useful next step is
to locate the inflection structure of $H_N$ exactly and determine
whether its Mellin integral restores monotonicity in the allocation
parameter. We make no untested assertion that $M=C_*$.

\paragraph{2. Quantify allocation monotonicity.}
The covariance formula proves a strict sign but does not give a sharp
modulus. Find useful two-sided estimates for
\[
 -\frac{d}{da}\{-g_{a,w-a}\}
 =-\operatorname{Cov}\!\left(J_w(V_a),
                          \log\frac{V_a}{1-V_a}\right),
\]
uniformly near $a=0$ or $a=w$, or as $w\to0$ and $w\to\infty$.
The endpoint beta laws concentrate at different rates and have
logarithmic score singularities. Deriving a full endpoint expansion
would turn the sharp inequalities into effective error estimates for
nearly extremal orders. It would also test whether higher derivative
signs hold on a restricted range; they do not follow from stochastic
ordering alone.

\paragraph{3. Extend the allocation principle to other arguments.}
For $z=e^{i\theta}$ the positive double integral still supplies an
explicit rational imaginary kernel, but its dependence on the
allocation variable can change sign and shape with $\theta$.
Characterize the angles for which an appropriate signed component is
monotone at fixed total order. This would connect the Gaussian theorem
to the incoming angular-zero program. In particular, the
below-threshold angular uniqueness conjecture from \cite{realorder}
remains open here: failure of a finite signed Hausdorff measure for
$a+b<1$ does not imply failure of the positive double-integral method.

\paragraph{4. Determine the analytic reach of the inverse expansion.}
Theorem~\ref{thm:inverse} gives an asymptotic expansion to every finite
order. Does a naturally grouped generalized-power series converge on a
nontrivial interval? The infinitely many exponents and their
multiplicative collisions make this a distinct question from a
finite-variable implicit-function theorem. A second target is a
uniform approximation joining the large-order branch to the
square-root behavior at the nondegenerate maximum. The latter behavior
follows locally from $C''(w_*)<0$; a global error-controlled transition
would require estimates beyond the local theorem.

\subsection{Integral distribution and cohomology}

\paragraph{5. Include ramified fixed grids.}
Remove the requirement that each primary factor be either entirely
fixed or fixed only at zero. The action $q=9,u=4$ gives the first
model. A resolution indexed by both denominator and fixed-point depth
should replace the present one-step fixed/moving split. The target is
an explicit integral Smith formula, not merely its reduction modulo
$\ell$. This distinction matters when higher powers of $\ell$ could
appear in more general symmetry groups.

\paragraph{6. Remove semisimplicity on the fixed factor.}
When $\ell\mid\varphi(m)$, character projectors cannot be obtained by
division by the unit-group order. The fixed Koszul complex still gives
a useful starting point, but generalized eigenspaces and nilpotent
operators replace scalar components. Determine whether the answer can
be stated using Fitting ideals of the operators $V_p-a_p$. Small exact
Smith computations should guide the statement before a general formula
is proposed.

\paragraph{7. Replace one jet parameter by several.}
For
\[
 B=\mathbb Z[t_1,\ldots,t_s]/(t_1^{L_1},\ldots,t_s^{L_s}),
\]
the contact ideal generated by the reduced nonconstant weights need
not be principal. The proof here already identifies the appropriate
Koszul homology. Its length, annihilator, and module structure are the
next invariants to compute. A single minimum valuation cannot capture
them. Complete-intersection cases and monomial contact ideals are
natural initial families with reproducible algebraic tests.

\paragraph{8. Compute composite cyclic and noncyclic symmetries.}
For a cyclic group of composite order, moving orbits can have several
stabilizers. For a noncyclic group, the periodic two-column argument
must be replaced by a fuller group resolution. Determine how the
prime-order formulas fit together under restriction and transfer.
An initial target is a product of two cyclic prime-order groups acting
on distinct primary factors. Classical universal norm distributions
\cite{cpc:Ouyang2002} provide an important comparison framework; the
novelty target should be an explicit finite integral formula for the
chosen weighted presentation.

\subsection{Global Lerch zero geometry}

\paragraph{9. Decide the persistent index-two monotonicity threshold.}
Theorem~\ref{lershape:thm:k2k3} proves that order three is the smallest
individual derivative order for which both positive zero branches
decrease throughout $0\le\rho\le1$. It does not prove that every
$k\ge3$ has this property. With the notation of
Theorem~\ref{lershape:thm:criterion}, the persistent statement is
equivalent to the endpoint inequalities
\[
 C_k(r_k)\le0\qquad\text{for all integers }k\ge3.
\]
This is a precise reduction of the remaining question to one sign per
order. A uniform large-$k$ expansion with a rigorous remainder,
combined with finitely many exact certificates, is a plausible route.
Such a proof should keep the zero-saturation threshold and the
branch-monotonicity threshold distinct.

\paragraph{10. Certify the order-two turning point itself.}
Existence and uniqueness of its nondegenerate minimum are already
proved here. Its approximate coordinates are close to the singular
endpoint $\rho=1$, which makes naive finite differences unreliable.
Produce a rational isolating rectangle for the simultaneous equations
$F_{2,2}^{(\rho)}(a)=0$ and
$\partial_\rho F_{2,2}^{(\rho)}(a)=0$, with a two-variable interval
Newton certificate. This would upgrade the present diagnostic
coordinates without changing the proved branch-shape theorem.

\paragraph{11. Extend the stationary-point bound to higher index.}
For general $n$, the Appell kernel bounds the number of positive
$a$-zeros by $n$. The resolvent theorem already controls the largest
one without assuming the bound is saturated. What is the optimal
number and type of stationary points on the other simple branches?
The index-two proof uses three coefficient-sign blocks. Higher index
suggests a finite combinatorial constraint, but a sign-change count
alone does not determine the type or global order of the stationary
points. A useful theorem would combine multiplicity at a stationary
parameter with the known sign of $F_a$ on that branch.

\subsection{Weight-seven identities and formal verification}

\paragraph{12. Enlarge the exact $S_6$ relation system strategically.}
The integer separator provides a direct test for any proposed new row
in the same ambient word space:
if its pairing is nonzero, it crosses the presently certified
obstruction. One route is to add convergent linear combinations of
Cayley words beyond the balanced seeds. Another is to allow
higher-depth coordinates and then eliminate them. A third is to seek
a functional identity before evaluation at $i$. Standard relations
and their possible incompleteness are discussed in \cite{zhao,racinet}.
The target remains the concrete $S_6$ identity, not a blanket assertion
that one selected relation family is complete.

\paragraph{13. Formalize the finite and analytic interfaces.}
The separator checker, Smith examples, and outward interval primitives
have small trusted mathematical interfaces. Formalizing these first
would provide useful infrastructure without immediately formalizing
the full special-function library. On the analytic side, the beta
covariance lemma and the strict normalized Mellin comparison are
reusable standalone lemmas. A good dependency order is positivity and
dominated differentiation, then the stochastic comparison, then the
special-function endpoint identities. This makes the formalization
serve several results rather than duplicate each proof separately.

These questions preserve the scope of the theorems already established.
The strongest immediate targets are the persistent index-two endpoint
sign, the ramified primary distribution calculation, and an $S_6$ row
with nonzero pairing against the supplied separator. Each has a precise
success criterion and a computational tool in the present package.
